The Districting Integrity Act represents the most complete available solution to partisan gerrymandering: federal legislation that eliminates partisan manipulation at its source by specifying a neutral algorithm whose output any person can verify. Its constitutional authority is grounded in Article I, \S4 and Fourteenth Amendment \S5. Its technical implementation is complete, available as the open-source \texttt{bisect} platform. Its historical precedent---the Apportionment Act of 1941---demonstrates that Congress has successfully resolved analogous disputes by delegating to mathematics.

The DIA outperforms all four alternative designs on the combined dimensions of constitutional durability, VRA compatibility, and technical implementability. It ranks below the criteria-only and court-remedy designs on political feasibility, but this disadvantage is a function of current political conditions rather than intrinsic design defects. A criteria-only bill with quantitative thresholds is a viable second-best option for Congress that lacks the votes for full algorithm specification.

The primary political barrier is incumbency self-interest. Members of Congress who benefit from safe districts are unlikely to vote for legislation that makes those districts competitive, regardless of the reform's merits. This barrier, analyzed in P.0, is real. But it is not insurmountable. The reform coalition must be built from members who are in competitive districts and who have less to lose from neutral maps, combined with members who represent geographically concentrated partisan minorities who are currently systematically underrepresented.

The timing window for federal action is 2025--2030. Redistricting for the 2030 census will occur in 2031. Federal legislation enacted by 2028 would give states two to three years to implement the \texttt{bisect} platform before the next redistricting cycle. Implementation is not technically complex---a single \texttt{bisect build} command produces the baseline plan in under 20 minutes---but states will need time to train personnel, conduct audits, and develop public engagement processes. Early enactment is essential.

If the DIA is enacted, the redistricting fights of 2031 will be a fraction of the battles of 2021. The algorithm will produce the baseline. VRA modifications will be documented and verifiable. Courts will review manifests, not cartographic expertise. The political energy that currently goes into gerrymandering will go somewhere else. That is the DIA's ultimate promise: not a perfect map, but a map whose imperfections are grounded in legal requirements rather than partisan strategy.
